% ============================================================================
%  CONCLUSIONES
% ============================================================================

\section{Conclusiones}

El desarrollo del m\'odulo de inscripci\'on a eventos deportivos en
la aplicaci\'on \nombreApp{} demuestra la integraci\'on exitosa de
m\'ultiples conocimientos trabajados durante la materia de
Aplicaciones M\'oviles I:

\begin{enumerate}
  \item \textbf{An\'alisis y dise\~no:} Se identific\'o la necesidad
    de un mecanismo de inscripci\'on a eventos, se defini\'o la
    funcionalidad, y se dise\~naron las interfaces manteniendo la
    consistencia visual del proyecto.

  \item \textbf{Implementaci\'on de interfaz:} Se crearon nuevos
    componentes (\texttt{EventRegistrationButton},
    \texttt{MyRegistrationsSection}) que se integran
    arm\'onicamente con la aplicaci\'on existente.

  \item \textbf{Navegaci\'on:} El m\'odulo se accede desde la
    pantalla de detalle del escenario (secci\'on de eventos) y
    desde la pantalla de perfil (secci\'on ``Mis inscripciones'').

  \item \textbf{L\'ogica de negocio:} Las operaciones de
    inscripci\'on, cancelaci\'on y consulta se implementaron
    mediante hooks personalizados de React Query.

  \item \textbf{Validaciones:} Se implementaron validaciones de
    autenticaci\'on, duplicados, integridad referencial y
    confirmaci\'on de acci\'on.

  \item \textbf{Manejo de estado:} React Query gestiona
    eficientemente el estado del servidor, la cach\'e y la
    invalidaci\'on selectiva.

  \item \textbf{Persistencia:} Toda la informaci\'on se almacena
    en Supabase (PostgreSQL) con pol\'iticas RLS granulares.

  \item \textbf{CRUD:} Se implementaron operaciones CREATE, READ
    y DELETE sobre la tabla \texttt{event\_registrations}.

  \item \textbf{Manejo de errores:} Se contemplaron errores de
    red, datos inv\'alidos, duplicados y usuarios no
    autenticados.

  \item \textbf{Integraci\'on:} El m\'odulo forma parte
    integral de la aplicaci\'on y no funciona como un proyecto
    independiente.
\end{enumerate}

El m\'odulo cumple con todos los requisitos m\'inimos establecidos
en la gu\'ia del examen integrador (PDF~7) y demuestra la
capacidad de resolver un desaf\'io real de desarrollo de software
m\'ovil utilizando las tecnolog\'ias trabajadas durante la materia.
